\chapter{绪论}

\section{研究背景及意义}

随着信息技术、工程教育改革、在线教学模式的不断发展，程序设计类课程对高校人才培养模式中的作用越来越显著。程序设计教学不仅要求学生掌握语法知识，还要求其拥有算法设计、问题抽象、调试分析和工程实现能力。如何在大规模教学场景中达成高效率、客观化且可重复的程序作业评价就成为了计算机类课程建设的关键问题之一，而在线编程评测平台的出现，为这一问题提供了较为有效的解决路径，在线编程评测平台借助自动编译、运行、测试用例比对、结果反馈等机制给程序设计教学予以了标准化支持，显示出较高的教学应用价值\cite{kurnia2001online,leal2003mooshak,ala2005survey}。

传统的在线编程评测平台大多只围绕提交即判分这一基本模式展开，虽然这样能够完成结果判定，但在教学反馈、资源弹性和平台运维等方面存在明显短板。一方面，单体式系统在高并发提交场景下容易受到资源瓶颈影响，另一方面，学生在面对编译错误、运行错误或答案错误时只能获得较为简短的判题结论，难以快速定位问题原因。对于以过程性学习和能力培养的教学场景而言，仅依赖传统的评测逻辑已难以满足更高层次的系统设计需求\cite{ala2005survey,keuning2018feedback,barczak2023automated}。

与此同时，云原生、容器化和微服务架构的快速发展为评测平台提供了更具可操作的技术路径。判题任务在教学场景里经常呈现出短时高峰、语言环境差异大、隔离要求严格等特征，这些特征与容器执行、服务拆分、弹性调度天然契合\cite{burns2016borg,balalaie2016microservices,difrancesco2019microservices,hezhenwei2020kubernetes}。将这些技术引用到系统中的价值并不仅限于部署到云上，还在于提升判题环境一致性、扩展弹性和长期维护可控性。

近年来，大语言模型和代码智能技术的发展又让在线评测平台具备了进一步提升反馈质量的可能。传统平台大多只能返回 \texttt{AC}、\texttt{WA}、\texttt{CE} 等简短状态，而实际应用场景中学生真正关心的是代码错在哪里、为什么会错以及应该怎么修改。在这方面，AI 技术就适合承担解释和提示的作用，但是不适合直接让 AI 进行判题或替代判题的主体\cite{keuning2018feedback,sarsa2022llm,jacobs2024feedback,wangfang2023ai,zhanghongzhuo2024genai}，如何在保证判题公平的前提下合理接入 AI 辅助反馈机制正是本文关注的重要问题之一。

基于上述背景，本文选择在线编程评测平台作为研究对象，并尝试在其中引入云原生执行环境与 AI 辅助反馈两类技术手段来解决两个较为实际的问题，如何将自动判题平台做得更稳定、更安全、更便于扩展以及如何在不影响判题公平性的前提下为学生补充更有用的解释性反馈。

\section{国内外研究现状}

综合上述工作可以看到，自动评测研究已经回答了如何保证判题正确与可重复，云原生研究提供了如何支撑弹性执行与资源约束，AI 教育研究则扩展了如何生成更具解释性的学习反馈。问题在于，这三类成果往往在不同系统中各自推进。面向评测业务的并发调度细节，常被云侧通用框架弱化，AI 反馈也常作为外挂能力存在，难以和判题流程共享上下文。本文的切入点正是缩小这三者之间的工程断层。

国外学者于自动判题、在线编程评测领域收获了诸多成果。Kurnia 等人讲述了 Online Judge 系统的设计运行状况，证明自动评测在大规模教学与训练场景具备可行性\cite{kurnia2001online}，Leal 和 Silva 所提出的 Mooshak 系统让竞赛与课程场景下的题目组织、提交流程、结果反馈机制更为规范\cite{leal2003mooshak}。Ala-Mutka 在自动评测方法综述当中，从功能正确性、效率、风格等多个方面对评估方法予以了比较全面的梳理\cite{ala2005survey}，Ihantola 等人在程序作业自动评估综述里剖析了自动反馈与教学场景的关系\cite{ihantola2011automated}，Keuning 等人又做了进一步说明。过程性、解释性反馈对学生学习编程有着不小的帮助，不过许多平台在反馈方面的程度仍有进步空间\cite{keuning2018feedback}。Paiva 等人在针对计算机科学作业的自动评测综述里观察到，效率、行为还有可读性等维度不断进入评估范围，容器化和静态分析技术助力评测平台的工程形态变得更完善\cite{paiva2022automated}。Messer 等人经过系统综述总结了自动评分与反馈工具的研究热点，表明评测维度、自动化程度、评估方法之间不均衡情况显著，公开数据集不足同样限制了不同工具之间的横向比较\cite{messer2024automated}。Barczak 等人在最新综述中归纳了自动评测平台的多条技术路线\cite{barczak2023automated}。

在云原生、容器化这一领域，Burns 等人针对 Borg 等大规模集群管理系统展开总结，其结果显示，容器编排、统一调度机制在应对短时突然出现的负载、实现资源的共享、多节点之间的协同工作这些方面，具备显著优势\cite{burns2016borg}，Balalaie、Di Francesco 等人围绕微服务架构、实证分析做了进一步阐述，他们指出，依据业务职责对模块进行拆分，对于独立部署、独立演化、故障隔离而言，有着积极的意义\cite{balalaie2016microservices,difrancesco2019microservices}，Docker 与 Kubernetes 的官方文档，还有 NIST 容器安全指南，从工程实践、安全规范这两个方面，为容器化的执行提供了相对完整的参考\cite{dockerdocs,kubernetesdocs,nist201716800190}。在 AI 辅助编程反馈这一方向上，Sarsa 等人探索了借助大语言模型来生成编程练习、代码解释的可行性\cite{sarsa2022llm}，Jacobs 等人对大语言模型自动反馈在入门编程课程里的实际效果做了评估\cite{jacobs2024feedback}。Finnie-Ansley 等人进行了关于 Codex 在入门编程教学中影响的研究，指出代码生成模型可能会改变学生练习和考核的形态，所以要重新审视学术诚信、能力培养目标\cite{finnieansley2022robots}，Denny 等人通过 Copilot 提示工程实验，发现自然语言与代码协同交互正在形成新的学习路径\cite{denny2023conversing}，Kazemitabaar 等人的实验研究表明，AI 代码生成器在提高新手提示效率的时候也可能会引发过度依赖\cite{kazemitabaar2023chi}。Raihan 等人通过系统文献综述梳理了大语言模型进入计算机教育的多条路径，包括课程建设、学习支持与风险治理等方向\cite{raihan2025llmcsed}，Koutcheme 等人证实开源大模型在生成反馈、作为评判者去评估反馈质量这两项任务方面有着实际潜力\cite{koutcheme2024evaluatinglm,koutcheme2024openfeedbackiticse}，Lubin 等人探讨了大语言模型在编程教学中的能力界限、使用规范\cite{lubin2025llm}。

国内研究起步相对略晚，但近年来在 OJ 平台、容器化教学环境与 AI 辅助教学等方面均有较快推进。Tran 等人在面向高校教学场景的在线评测系统 UOJ 中，对题目管理、提交流程与判题机制做了较为完整的工程实现\cite{tran2013uoj}，何振卫等人围绕 Kubernetes 的资源调度展开研究，对容器编排在教学与实验环境中的应用进行了讨论\cite{hezhenwei2020kubernetes}，陈娟等人针对在线编程教学环境中的反馈机制开展研究，强调过程性反馈对学习成效的影响\cite{chenjuan2020online}。在 AI 辅助教学方面，王芳等人结合教学实践讨论了人工智能在程序设计课程中的应用模式\cite{wangfang2023ai}，张洪卓等人则面向生成式 AI 在高校教学中的影响给出了较为系统的分析\cite{zhanghongzhuo2024genai}，Alemu 等人提出的无服务器化自动评测方案也代表了在轻量化、低成本部署方向上的一类探索\cite{alemu2023serverless}。

图~\ref{fig:ch1-research-map}给出了对应脉络与交汇关系。

\begin{figure}[htbp]
  \centering
  \resizebox{0.88\textwidth}{!}{% 第1章：三条研究脉络与本文交汇示意
\begin{tikzpicture}[
  font=\small,
  dmellipse/.style={draw, ellipse, minimum width=3.0cm, minimum height=1.1cm, align=center, inner sep=2pt},
  focus/.style={draw, very thick, rounded corners=4pt, fill=yellow!12, inner sep=8pt, align=center}
]
  \node[dmellipse, fill=cyan!10] (A) {在线编程评测\\与自动判题};
  \node[dmellipse, fill=lime!15, right=1.8cm of A] (B) {云原生与\\弹性执行};
  \node[dmellipse, fill=magenta!12, right=1.8cm of B] (C) {AI 辅助\\编程反馈};
  \node[focus, below=1.5cm of B] (X) {\textbf{本文}\\一体化在线评测平台\\（判题主链 + 容器执行 + 解释性反馈）};
  \draw[-{Stealth}, thick] (A.south) -- (X.north west);
  \draw[-{Stealth}, thick] (B.south) -- (X.north);
  \draw[-{Stealth}, thick] (C.south) -- (X.north east);
\end{tikzpicture}
}
  \caption{在线判题研究、云原生工程与 AI 教育应用三条脉络及其交汇关系示意}
  \label{fig:ch1-research-map}
\end{figure}

\section{本文主要工作}

全文依次经历文献与背景梳理、需求分析、总体设计、工程实现、测试验证与总结展望：第二章归纳判题、容器与异步、反馈分层及 AI 使用边界等共性技术基础\cite{kurnia2001online,leal2003mooshak,ala2005survey,keuning2018feedback,barczak2023automated,burns2016borg,balalaie2016microservices,difrancesco2019microservices,sarsa2022llm,jacobs2024feedback,chenjuan2020online,hezhenwei2020kubernetes,wangfang2023ai,zhanghongzhuo2024genai,tran2013uoj,alemu2023serverless,ihantola2011automated,lubin2025llm}；第三章面向教学场景给出功能与非功能需求；第四章给出总体架构、异步判题与数据模型、工作区与安全边界等设计要点；第五章在既定技术栈上实现核心链路并说明关键模块与运行证据；第六章通过功能、判题、安全与交互等用例给出验证结论。

本文目标是在高校程序设计教学场景下，形成一套可稳定演示的在线评测核心链路（题库、提交、自动判题、结果与过程数据、可选 AI 辅助），论述范围不以「大而全教务系统」为中心。贯穿各章的主线是：客观判题与 AI 解释分层、容器化执行与异步调度、提交—判题—用例分层留存，以及题面与编码同屏的一体化工作区；技术概念、方案细节、工程路径与测试证据分别集中在第二、四、五、六章，本章不展开实现级描述。

全文章节安排如下：第一章绪论；第二章相关技术基础；第三章需求分析；第四章总体设计；第五章核心实现；第六章系统验证；第七章总结与展望。
